\section*{Capítulo \ref{ch:g4:fractions} --- Primeiras frações}

\begin{solution}{exo:g4:fractions:1}
Pizza: sobra $\frac56$. Chocolate: $\frac38$ comido. Bandeira:
$\frac13$ colorida.
\end{solution}

\begin{solution}{exo:g4:fractions:2}
$\frac35$ pinta três partes e $\frac25$ pinta só duas: $\frac35$ é a
pintura maior.
\end{solution}

\begin{solution}{exo:g4:fractions:3}
\omterm{def:g3:division:half}{Metade} de $18$: $9$. Quarto de $20$: $5$. Terço de $21$: $7$.
$\frac34$ de $20$: três quartos, $3 \times 5 = 15$.
\end{solution}

\begin{solution}{exo:g4:fractions:4}
$\frac13$: um passo; $\frac33$: em cima do $1$; $\frac53$: dois
passos depois do $1$. A \omterm{def:g4:fractions:def}{fração} que cai exatamente sobre o $2$ é
$\frac63$.
\end{solution}

\begin{solution}{exo:g4:fractions:5}
$\frac38 < \frac68$; \quad $\frac55 = 1$; \quad $\frac74 > 1$;
\quad $\frac12 = \frac24$.
\end{solution}

\begin{solution}{exo:g4:fractions:6}
$\frac26 + \frac36 = \frac56$; \quad
$\frac58 - \frac28 = \frac38$; \quad
$\frac14 + \frac24 = \frac34$; \quad
$1 = \frac33$, então $1 - \frac13 = \frac23$.
\end{solution}

\begin{solution}{exo:g4:fractions:7}
$\frac54 = 1 + \frac14$; \quad
$\frac73 = 2 + \frac13$; \quad
$\frac92 = 4 + \frac12$.
\end{solution}

\begin{solution}{exo:g4:fractions:8}
Meia hora: $30$ min. Um quarto: $15$ min. Três quartos: $45$ min. Um
terço: $20$ min.
\end{solution}

\begin{solution}{exo:g4:fractions:9}
$\frac14$ de $24$ são $6$ alunos com óculos; $24 - 6 = 18$ sem.
\end{solution}

\begin{solution}{exo:g4:fractions:10}
Juntos: $\frac38 + \frac48 = \frac78$. Sobrou: $\frac18$. Samuel
comeu mais ($4$ oitavos contra $3$).
\end{solution}

\begin{solution}{exo:g4:fractions:11}
É possível mesmo: \omterm{def:g3:division:half}{metade} de uma pizza pequena pode ser menos comida
do que um quarto de uma pizza gigante. As frações comparam partes
\emph{do mesmo inteiro}; antes de comparar $\frac12$ e $\frac14$ ``de
alguma coisa'', é preciso saber que as duas coisas são iguais.
\end{solution}
